\documentclass[11pt]{article}
\usepackage[top=.78in,bottom=.76in,left=.83in,right=.83in]{geometry}
\usepackage{amsmath,amssymb,amsthm,booktabs,microtype,xcolor}
\usepackage[colorlinks=true,urlcolor=blue!45!black]{hyperref}
\newtheorem{theorem}{Theorem}
\newcommand{\F}{\mathbb F}
\newcommand{\Q}{\mathbb Q}
\begin{document}
\begin{center}
{\LARGE A sharper five-step graph-tower bound}\\[5pt]
{\large One depth-five direction, and exact maxima at larger primes}\\[6pt]
{\small 27 November 2026}
\end{center}
\begin{abstract}
We lower the general five-step ceiling from $4(p+1)+18$ to
$4(p+1)+16$ outside $p=11,31$. Precisely one depth-five direction
gives excess at most $14$, attained for infinitely many primes.
A finite incidence search makes this a statement about the
unrestricted optimum at $1031,1511,1571,2351$, and also separates
simple and unrestricted optima at $761$ and $2281$.
A separate table supplies fifteen exact benchmarks through $181$.
\end{abstract}
\section{Bounds from the quartic and sextic moments}
Let $p>10$ be prime and $S=(s_0,\ldots,s_4)$ a spanning list of
residues in $\F_p^2$. Use the notation of the
\href{https://github.com/VladimirReshetnikov/ProveIt/blob/f6b5955e13bcefb2760480e7b823eba55ed2383e/Analysis/GeneralizedRiemannHypothesis/Article/sections/gra.tex}{merged ProveIt article}:
$M_j(u)=\sum_i(u\cdot s_i)^{2j}$, $r(L)$ the first nonzero
moment index in projective direction $L$, and $A=2\sum_Lr(L)$.
For the spanning-tree count $\kappa_n$ at layer $n$,
$v_p(\kappa_n)=A(p^n-1)/(p-1)-2n$.
Five signed-distinct nonzero step residues give simple graphs at every layer.

\begin{theorem}
Outside $p=11,31$, $A\le4(p+1)+16$. If exactly one direction has
depth five, then $A\le4(p+1)+14$. The latter bound is attained
by simple towers for infinitely many primes.
\end{theorem}

Newton's identities give $r(L)\le5$. If $M_1\ne0$, the quadratic
has at most two projective roots, so $A\le2(p+1)+16$, below both
claimed ceilings. If $M_1=0$, then $M_2$ is a nonzero quartic
by the article's first-moment lemma, and the excess above $4(p+1)$
is $2\sum(r-2)$ over its at most four projective roots.
With no depth-five direction this gives excess at most $16$.
With at least two, the article's signed-permutation certificate
(\texttt{gra:thm:fivestep}) gives excess at most $16$ outside
$11,31$, and at most $12$ outside $11,31,101,131$.
It remains to handle a single depth-five direction.

The squared projections in that direction solve $T^5-c=0$
with $c\ne0$. They give a coset of fifth roots of unity among
nonzero squares, so $p\equiv1\bmod10$.
Scale the row, reorder the steps and choose column signs so that
the row is $a=(1,z,z^2,z^3,z^4)$ for $z$ of order five.
The componentwise powers $a^j$, $0\le j<5$, form a basis.
A second row $b$ satisfies $a\cdot b=b\cdot b=0$.
The first condition removes the $a^4$ coefficient; subtracting
a multiple of $a$ removes its $a$ coefficient. Writing
$b=B a^0+C a^2+D a^3$, the second condition gives $B^2+2CD=0$.
If $C=0$, then $b$ is proportional to $a^3$. Otherwise scale
$b$ to obtain
\[
                        b=t a^0+a^2-\tfrac12t^2 a^3.
\]
Both $b=a^3$ and $b=a^2$ have two depth-five directions by
geometric sums, so only $t\ne0$ is relevant.

\section{Exact polynomial certificates}
In the finite directions $b+v a$, set
$U=2t+2vX+2X^2-t^2X^3$ in the algebra with $X^5=1$.
Let $f$ and $g$ be the constant coefficients of $U^4$ and $U^6$;
they differ from the moments only by nonzero scalars since the
sum of all fifth roots to a power is $0$ or $5$.
Put $s=t^5$. Exact expansion and division give
\[
\begin{split}
 f&=64v^3-96t^3v^2-8t(s-24)v-56t^4,\\
 g\bmod f&=3R,\\
 R&=t(128+352s+367s^2)
       +t^3(-768-1122s+53s^2)v+(320+670s^2)v^2.
\end{split}
\]
The gcd of the three pairwise resultants of the coefficient
quadratics in $s$ is $2^9 5^4 11^2$, so $R$ is not identically
zero in $v$ in this case.
For two distinct common roots, $R$ must be quadratic and divide
$f$. Multiplying $f\bmod R$ by the square of its leading
coefficient (that of $R$) gives $-32t^4P(s)-32tQ(s)v$, where
\[
\begin{aligned}
 P&=252928+489984s+1480832s^2+78716s^3+9003s^4,\\
 Q&=-532480-191488s-4536064s^2-283272s^3
                         +291464s^4+77s^5.
\end{aligned}
\]
The absolute resultant of $P,Q$ is
$2^{42}5^{24}11^{10}13^2 31^2\cdot761\cdot1297$.
The only possible prime $1\bmod10$ besides $11,31$ is $761$.
There a direct check over all nonzero $t$ finds no common root
of $f,R$. Consequently at most one of the three finite
exceptional directions can have depth four; the excess is at
most $6+4+2+2=14$. This finishes the upper bounds.

For attainment, the exact identities
\[
 \operatorname{Res}_v(f,R)=-1024t^3C_0(t^5),\qquad
 \operatorname{disc}_v(f)=65536t^3D_0(t^5)
\]
are verified in the checker with explicit polynomials $C_0,D_0$.
Here $C_0$ is nonconstant with nonzero constant term, and
$C_0,D_0$ are coprime by reduction modulo three with degrees
unchanged. Choose $t\ne0$ in characteristic zero with
$C_0(t^5)=0$; $f$ then has three distinct roots and exactly one
is common with $R$, by the preceding remainders.
Take a Galois number field containing $t$, the three roots, and
a primitive fifth root. Infinitely many primes split completely
there by the Chebotarev density theorem. Discard finitely many
primes where these identities, nonvanishings, and distinctness
fail. The finite roots give two depth-three directions and one
of depth at least four. After excluding $11,31,101,131$, the
two-depth-five certificate allows excess at most $12$ if
another depth-five direction exists. The observed pattern would
then already give at least $16$, so its last finite depth is
exactly four. This proves infinite attainment, with simple
steps since the first row has five distinct signed classes.

\section{Exact maxima at six larger primes}
Here is a second search which does not enumerate all step lists.
If the quartic has four directions of depth at least four,
the row plane contains four points of
$C:\sum_i x_i^2=\sum_i x_i^4=\sum_i x_i^6=0$ in $\mathbb P^4$.
Conversely the intersections with $C$ are exactly those
directions. To enumerate the points, use a permutation to put
a nonzero coordinate first, and scale it to $1$. For squared
second and third coordinates $a,b$, the squared last two
coordinates have sum $\sigma=-1-a-b$ and product
$\tau=(\sigma^2+1+a^2+b^2)/2$. The third power sum is zero exactly
when $\sigma^3-3\sigma\tau+1+a^3+b^3=0$. Test this identity and then
the quadratic $T^2-\sigma T+\tau$ for two square roots in the sense
that its two roots must both be squares (zero is allowed).
This exhausts the points after restoring permutations and signs.
For each first point up to those symmetries, pair it with every
distinct orthogonal second point. The span is isotropic, and
evaluating its remaining projective directions gives the full
count of points of $C$ on that plane.

At $1031,1511,1571,2351$ no such plane has four points of $C$,
so the zero-depth-five case gives at most $14$; the one-depth-five
case attains $14$ as checked below. At $2281$ every plane with four
points of $C$ has a zero step, so the same holds when all residues
are nonzero. At $761$ every plane with nonzero steps has at most two
points of $C$. A plane without a depth-five direction can then have
at most excess $12$. With exactly one depth-five direction, the
exceptional modular check for $761$ gives no other depth-four
direction, again yielding at most $12$. More than one gives at
most $12$ by the previous theorem. The simple family with rows
$a,a^2$ attains $12$. At $761$ and $2281$, which are $1\bmod8$,
append a zero to an optimal four-step list to attain $16$.
\begin{center}
\begin{tabular}{rrrrr}
\toprule
$p$ & $z$ & $t$ & Zero allowed & Nonzero and simple\\
\midrule
761  & --  & --   & 3064 & 3060\\
1031 & 264 & 381  & 4142 & 4142\\
1511 & 534 & 593  & 6062 & 6062\\
1571 & 382 & 213  & 6302 & 6302\\
2281 & 342 & 170  & 9144 & 9142\\
2351 & 207 & 1124 & 9422 & 9422\\
\bottomrule
\end{tabular}
\end{center}
For the five non-dash rows the table gives the parameters of the
simple witness $(a,b)$ above; the exact moment checks are in the
certificate. For instance, at $1031$ the exact maximal valuation
is $4142(1031^n-1)/1030-2n$ for every $n$.

\section{Fifteen additional exact benchmarks}
The complementary-key search for $M_1=0$ uses two steps $e_1,e_2$,
and three further signed representatives. For each sorted pair
of further steps, determine the remaining signed class from the
sum of the three keys $(x^2,xy,y^2)$, which must be $(-1,0,-1)$.
The map is injective on sign classes, including zero.
Check every surviving triple's moments. At each prime below,
all maxima exceed the $M_1\ne0$ bound. Simple witnesses and
normalization counts are supplied in the output.
\begin{center}
\small
\begin{tabular}{rrr @{\qquad} rrr}
\toprule
$p$ & Zero allowed & Nonzero / simple &$p$ & Zero allowed & Nonzero / simple\\
\midrule
103&424&424&149&608&608\\
107&440&440&151&620&620\\
109&448&448&157&640&640\\
113&472&466&163&664&664\\
127&520&520&167&680&680\\
137&568&562&173&708&708\\
139&568&568&179&728&728\\
&&&181&740&740\\
\bottomrule
\end{tabular}
\end{center}
The results are unconditional, and sharpen but do not solve the
five-step maximum problem at arbitrary primes.
\end{document}
